\section*{Limitations}
The primary limitation is the label source. PubMed-PICO and PubMed RCT labels are derived from abstract structure or annotation conventions, and a section label is not always equivalent to sentence-level semantics. Some errors identified in the analysis likely reflect ambiguous or proxy labels rather than purely model failures.

The external validation corpus is public and held out, but it remains in the biomedical abstract domain. It is not a manually adjudicated external clinical setting, and it does not prove deployment readiness. The NICTA-PIBOSO corpus was investigated as a manual external validation candidate but was not automatically available in usable form in this environment.

The neural experiments were intentionally CPU-limited. Models used small embeddings, short maximum sequence lengths, and a stratified cap on neural training examples. Larger recurrent models, pretrained biomedical Transformers, or GPU-scale fine-tuning might perform differently, but they were outside the CPU-only scope of this project.

Calibration uncertainty is another limitation. The project reports calibrated Brier score, expected calibration error, and reliability curves, but it did not retain pre-calibration Brier/ECE for all models and did not compute uncertainty intervals for calibration metrics. No demographic fairness analysis was possible because the public abstract corpora do not include patient-level demographic metadata.
